% 【本文件写什么】impl 实现层：block-cache
% - 定位：写穿 LRU 块缓存装饰器。
% - crate 路径：os/components/wateros-driver/driver-block/block-impl/impl-block-cache/src/
% - 如何实现对应 api-v0 trait：关键类型、状态机、数据结构
% - 算法或硬件交互流程，可用伪代码/流程图
% - 本 impl 启用的 Cargo [features] 与 arch feature，impl-riscv64 等
% - 与其它 impl 变体的差异、切换 feature 时的影响
% - 性能/内存假设与已知限制
% 【事实来源】
% - 上述 crate 源码与 Cargo.toml
% - os/feature-tree.txt 中指向本 impl 的 feature 链

\subsubsection{impl-block-cache}

\paragraph{启用链。}根\code{wateros}默认\code{qemu-riscv64-opensbi}同时打开\code{driver/impl-block-cache}，向RISC-V与LoongArch平台impl传递\code{block-cache} feature。注册virtio-blk时平台调用\code{BlockCacheManager::wrap}而非直接\code{Arc::new(Mutex::new(dev))}。

\paragraph{写穿LRU策略。}\code{CachingBlockDevice}包装任意\code{Box<dyn BlockDevice + Send>}，对上仍实现\code{BlockDevice}。\code{capacity\_blocks}来自\code{wateros-base-config::BLOCK\_CACHE\_CAPACITY\_BLOCKS}；为0时退化为透传。\code{read\_blocks}：命中\code{BTreeMap<Lba, slot\_idx>}则避免访问\code{inner}；连续未命中区间合并为单次\code{inner.read\_blocks}后逐块\code{cache\_put}。\code{write\_blocks}：先写穿\code{inner}，再更新或插入缓存行。

\paragraph{淘汰与恢复。}槽位由\code{free}栈与\code{lru}双端队列管理；\code{evict\_lru\_slot}弹出最久未用槽并清除\code{map}条目。若淘汰遇到不变式破坏，\code{reset\_cache\_invariant}清空\code{map}/\code{lru}并重建\code{free}。\code{flush}在写穿语义下为no-op，保留供将来write-back或测试钩子。

\begin{rustcode}
pub struct CachingBlockDevice {
    inner: Box<dyn BlockDevice + Send>,
    block_size: usize,
    capacity: usize,
    map: BTreeMap<Lba, usize>,
    slots: Vec<Slot>,
    free: Vec<usize>,
    lru: VecDeque<usize>,
}
\end{rustcode}

\code{BlockCacheManager::default\_config}读取\code{base-config}容量；\code{wrap}返回可直接\code{register\_block\_device}的\code{SharedBlockDevice}。\code{flush\_all}当前恒\code{Ok(())}。单元测试覆盖同块重复读仅一次底层读连续未命中合并读写后读命中缓存等路径。
